\documentclass[11pt,a4paper]{article}
\usepackage[margin=2.5cm]{geometry}
\usepackage{amsmath,amssymb}
\usepackage{graphicx}
\usepackage{tcolorbox}
\usepackage{listings}
\usepackage{xcolor}
\usepackage{hyperref}
\usepackage{tikz}
\usetikzlibrary{shapes,arrows,positioning}

% Color definitions
\definecolor{mlblue}{RGB}{0,102,204}
\definecolor{mlpurple}{RGB}{51,51,178}
\definecolor{mlorange}{RGB}{255,127,14}
\definecolor{mlgreen}{RGB}{44,160,44}

% Key concept box
\newtcolorbox{keybox}[1]{
  colback=mlblue!5!white,
  colframe=mlblue,
  fonttitle=\bfseries,
  title=#1
}

% Example box
\newtcolorbox{examplebox}[1]{
  colback=mlgreen!5!white,
  colframe=mlgreen,
  fonttitle=\bfseries,
  title=#1
}

% Security box
\newtcolorbox{securitybox}[1]{
  colback=mlorange!5!white,
  colframe=mlorange,
  fonttitle=\bfseries,
  title=#1
}

% Code listing style
\lstset{
  basicstyle=\ttfamily\small,
  backgroundcolor=\color{gray!10},
  frame=single,
  numbers=left,
  numberstyle=\tiny\color{gray},
  breaklines=true,
  captionpos=b
}

\title{\textbf{Lecture Notes: Lesson 2}\\Cryptography and Security}
\author{Cryptocurrency Course}
\date{}

\begin{document}

\maketitle

\tableofcontents
\newpage

\section{Introduction to Cryptography in Blockchain}

Cryptography is the cornerstone of blockchain security. It provides:
\begin{itemize}
  \item \textbf{Data Integrity}: Ensuring data hasn't been tampered with
  \item \textbf{Authentication}: Proving identity without revealing secrets
  \item \textbf{Non-repudiation}: Preventing denial of actions taken
  \item \textbf{Confidentiality}: Protecting sensitive information (in some blockchains)
\end{itemize}

\begin{keybox}{Core Cryptographic Primitives}
Blockchain systems rely on two fundamental cryptographic tools:
\begin{enumerate}
  \item \textbf{Hash Functions}: For data integrity and chain linkage (SHA-256)
  \item \textbf{Digital Signatures}: For authentication and authorization (ECDSA)
\end{enumerate}
\end{keybox}

\section{SHA-256: Secure Hash Algorithm}

\subsection{Algorithm Structure}

SHA-256 belongs to the SHA-2 family designed by the NSA and published by NIST in 2001. It produces a 256-bit (32-byte) hash from arbitrary-length input.

\subsubsection{Internal Architecture}

The algorithm operates on 512-bit message blocks through:

\begin{enumerate}
  \item \textbf{Message Padding}:
  \begin{itemize}
    \item Append bit '1' to message
    \item Append '0' bits until length $\equiv 448 \pmod{512}$
    \item Append original message length as 64-bit integer
  \end{itemize}

  \item \textbf{Initialize Hash Values}: Eight 32-bit registers (H0-H7)
  \[
  \begin{aligned}
  H_0 &= \text{0x6a09e667} \quad H_1 = \text{0xbb67ae85} \\
  H_2 &= \text{0x3c6ef372} \quad H_3 = \text{0xa54ff53a} \\
  H_4 &= \text{0x510e527f} \quad H_5 = \text{0x9b05688c} \\
  H_6 &= \text{0x1f83d9ab} \quad H_7 = \text{0x5be0cd19}
  \end{aligned}
  \]
  These are the fractional parts of square roots of first 8 primes.

  \item \textbf{Message Schedule}: Expand 512-bit block to 64 words ($W_0 \ldots W_{63}$)

  \item \textbf{Compression Function}: 64 rounds of operations using:
  \begin{itemize}
    \item Logical functions: Ch, Maj, $\Sigma_0$, $\Sigma_1$, $\sigma_0$, $\sigma_1$
    \item Constants $K_t$ (fractional parts of cube roots of first 64 primes)
    \item Modular addition and bitwise operations
  \end{itemize}

  \item \textbf{Finalization}: Add compressed values to current hash values
\end{enumerate}

\subsection{Key Operations}

\subsubsection{Logical Functions}

\begin{align*}
\text{Ch}(x, y, z) &= (x \land y) \oplus (\neg x \land z) \\
\text{Maj}(x, y, z) &= (x \land y) \oplus (x \land z) \oplus (y \land z) \\
\Sigma_0(x) &= \text{ROTR}^2(x) \oplus \text{ROTR}^{13}(x) \oplus \text{ROTR}^{22}(x) \\
\Sigma_1(x) &= \text{ROTR}^6(x) \oplus \text{ROTR}^{11}(x) \oplus \text{ROTR}^{25}(x) \\
\sigma_0(x) &= \text{ROTR}^7(x) \oplus \text{ROTR}^{18}(x) \oplus \text{SHR}^3(x) \\
\sigma_1(x) &= \text{ROTR}^{17}(x) \oplus \text{ROTR}^{19}(x) \oplus \text{SHR}^{10}(x)
\end{align*}

where ROTR is rotate right and SHR is shift right.

\subsection{Avalanche Effect Demonstration}

The avalanche effect ensures that small input changes cascade through the hash, producing drastically different outputs.

\begin{examplebox}{Avalanche Effect in SHA-256}
\begin{lstlisting}[language=Python]
import hashlib

def demonstrate_avalanche():
    """Show how 1-bit change affects entire hash"""
    messages = [
        "Hello World",
        "Hello World!",  # Added 1 character
        "hello World",   # Changed 1 bit (H->h)
    ]

    for msg in messages:
        hash_val = hashlib.sha256(msg.encode()).hexdigest()
        print(f"Input:  '{msg}'")
        print(f"SHA256: {hash_val}")
        print()

demonstrate_avalanche()
\end{lstlisting}

Output:
\begin{verbatim}
Input:  'Hello World'
SHA256: a591a6d40bf420404a011733cfb7b190d62c65bf0bcda32b57b277d9ad9f146e

Input:  'Hello World!'
SHA256: 7f83b1657ff1fc53b92dc18148a1d65dfc2d4b1fa3d677284addd200126d9069

Input:  'hello World'
SHA256: 5eb63bbbe01eeed093cb22bb8f5acdc3299ae8ad6e45a5b9b8f0a35a33c15f95
\end{verbatim}

\textbf{Analysis}: Changing a single bit flips approximately 50\% of the output bits, demonstrating excellent avalanche properties.
\end{examplebox}

\subsection{Security Strength}

\begin{keybox}{Computational Security}
SHA-256's 256-bit output provides:
\begin{itemize}
  \item \textbf{Pre-image Resistance}: $2^{256}$ operations ($\approx 10^{77}$)
  \item \textbf{Collision Resistance}: $2^{128}$ operations (Birthday paradox)
  \item \textbf{Current Status}: No practical attacks known
\end{itemize}

For comparison: There are only $\approx 10^{80}$ atoms in the observable universe.
\end{keybox}

\section{Public Key Cryptography}

\subsection{Asymmetric Encryption Fundamentals}

Unlike symmetric encryption (same key for encryption/decryption), public key cryptography uses a key pair:

\begin{itemize}
  \item \textbf{Private Key} ($sk$): Secret, used for signing
  \item \textbf{Public Key} ($pk$): Shared publicly, used for verification
\end{itemize}

\textbf{Mathematical Property}:
\[
\text{Verify}_{pk}(\text{Sign}_{sk}(m)) = \text{True}
\]

but knowing $pk$ doesn't reveal $sk$ (computational one-way function).

\subsection{Elliptic Curve Cryptography (ECC)}

Bitcoin and Ethereum use ECC, specifically the \texttt{secp256k1} curve, for digital signatures.

\subsubsection{Elliptic Curve Definition}

An elliptic curve over finite field $\mathbb{F}_p$ is defined by:
\[
y^2 \equiv x^3 + ax + b \pmod{p}
\]

For \texttt{secp256k1}:
\begin{itemize}
  \item $a = 0$, $b = 7$
  \item $p = 2^{256} - 2^{32} - 2^9 - 2^8 - 2^7 - 2^6 - 2^4 - 1$
  \item Generator point $G$
  \item Order $n$ (prime number of points)
\end{itemize}

Equation simplifies to:
\[
y^2 \equiv x^3 + 7 \pmod{p}
\]

\subsubsection{Point Addition}

The curve has a group structure with point addition operation:

\textbf{Geometric Interpretation}: To add points $P$ and $Q$:
\begin{enumerate}
  \item Draw line through $P$ and $Q$
  \item Find third intersection point $R'$
  \item Reflect $R'$ across x-axis to get $R = P + Q$
\end{enumerate}

\textbf{Algebraic Formula}: For $P = (x_1, y_1)$ and $Q = (x_2, y_2)$:
\begin{align*}
\lambda &= \begin{cases}
\frac{y_2 - y_1}{x_2 - x_1} & \text{if } P \neq Q \\
\frac{3x_1^2 + a}{2y_1} & \text{if } P = Q \text{ (point doubling)}
\end{cases} \\
x_3 &= \lambda^2 - x_1 - x_2 \\
y_3 &= \lambda(x_1 - x_3) - y_1
\end{align*}

\subsubsection{Scalar Multiplication}

Key operation: $Q = k \cdot G$ (add $G$ to itself $k$ times)

\textbf{Efficient Computation}: Use double-and-add algorithm in $O(\log k)$ time.

\begin{keybox}{Elliptic Curve Discrete Logarithm Problem (ECDLP)}
Given $Q = k \cdot G$, it's computationally infeasible to find $k$ (the private key).

This is the foundation of ECC security. Best known attacks require $O(\sqrt{n})$ time, where $n \approx 2^{256}$ for secp256k1.
\end{keybox}

\section{Digital Signatures with ECDSA}

\subsection{Elliptic Curve Digital Signature Algorithm}

ECDSA provides authentication and non-repudiation for blockchain transactions.

\subsubsection{Key Generation}

\begin{enumerate}
  \item Choose random private key: $sk \in [1, n-1]$
  \item Compute public key: $pk = sk \cdot G$
\end{enumerate}

\begin{examplebox}{Key Generation Example}
\begin{lstlisting}[language=Python]
from ecdsa import SigningKey, SECP256k1
import hashlib

# Generate private key
private_key = SigningKey.generate(curve=SECP256k1)

# Derive public key
public_key = private_key.get_verifying_key()

# Export as hex
sk_hex = private_key.to_string().hex()
pk_hex = public_key.to_string().hex()

print(f"Private Key: {sk_hex}")
print(f"Public Key:  {pk_hex}")

# Generate Bitcoin-style address
pk_hash = hashlib.sha256(public_key.to_string()).digest()
address_hash = hashlib.new('ripemd160', pk_hash).hexdigest()
print(f"Address (simplified): {address_hash}")
\end{lstlisting}
\end{examplebox}

\subsubsection{Signature Generation}

To sign message $m$ with private key $sk$:

\begin{enumerate}
  \item Hash the message: $e = H(m)$ (where $H$ is SHA-256)
  \item Choose random nonce: $k \in [1, n-1]$
  \item Compute point: $P = k \cdot G = (x, y)$
  \item Calculate $r = x \bmod n$
  \item Calculate $s = k^{-1}(e + r \cdot sk) \bmod n$
  \item Signature is pair: $(r, s)$
\end{enumerate}

\begin{securitybox}{Nonce Security Critical}
The nonce $k$ must be:
\begin{itemize}
  \item \textbf{Random}: Predictable $k$ leaks private key
  \item \textbf{Unique}: Reusing $k$ for different messages leaks private key
  \item \textbf{Secret}: Revealing $k$ leaks private key
\end{itemize}

Historical disaster: PlayStation 3 used constant $k$, leading to complete system compromise in 2010.
\end{securitybox}

\subsubsection{Signature Verification}

Given message $m$, signature $(r, s)$, and public key $pk$:

\begin{enumerate}
  \item Hash the message: $e = H(m)$
  \item Calculate: $w = s^{-1} \bmod n$
  \item Calculate: $u_1 = e \cdot w \bmod n$
  \item Calculate: $u_2 = r \cdot w \bmod n$
  \item Calculate point: $P = u_1 \cdot G + u_2 \cdot pk = (x, y)$
  \item Verify: $r \equiv x \pmod{n}$
\end{enumerate}

If $r \equiv x \pmod{n}$, signature is valid.

\subsection{Mathematical Proof of Correctness}

\begin{keybox}{Why ECDSA Works}
During signing:
\[
s = k^{-1}(e + r \cdot sk) \implies k = s^{-1}(e + r \cdot sk)
\]

During verification:
\begin{align*}
P &= u_1 \cdot G + u_2 \cdot pk \\
  &= (e \cdot s^{-1}) \cdot G + (r \cdot s^{-1}) \cdot (sk \cdot G) \\
  &= s^{-1}(e + r \cdot sk) \cdot G \\
  &= k \cdot G \quad \text{(same point computed during signing)}
\end{align*}

Therefore, $x$-coordinate matches $r$, confirming authenticity.
\end{keybox}

\section{Wallet Security}

\subsection{Wallet Types}

\subsubsection{Hot Wallets}

Connected to the internet:
\begin{itemize}
  \item \textbf{Web Wallets}: Browser-based (e.g., MetaMask)
  \item \textbf{Mobile Wallets}: Smartphone apps
  \item \textbf{Desktop Wallets}: Computer software
\end{itemize}

\textbf{Pros}: Convenient for frequent transactions \\
\textbf{Cons}: Vulnerable to malware, phishing, server breaches

\subsubsection{Cold Wallets}

Offline storage:
\begin{itemize}
  \item \textbf{Hardware Wallets}: Dedicated devices (Ledger, Trezor)
  \item \textbf{Paper Wallets}: Printed private keys
  \item \textbf{Steel Wallets}: Engraved metal plates
\end{itemize}

\textbf{Pros}: Maximum security for long-term storage \\
\textbf{Cons}: Less convenient for frequent use

\subsection{Hierarchical Deterministic Wallets (HD Wallets)}

HD wallets (BIP-32 standard) generate multiple addresses from a single seed.

\subsubsection{Derivation Path}

\begin{verbatim}
Master Seed (128-256 bits)
    |
    v
Master Private Key (m)
    |
    +--> Child Key (m/0)
    |       |
    |       +--> Grandchild (m/0/0)
    |
    +--> Child Key (m/1)
            |
            +--> Grandchild (m/1/0)
\end{verbatim}

\textbf{Standard Paths}:
\begin{itemize}
  \item Bitcoin: \texttt{m/44'/0'/0'/0/i}
  \item Ethereum: \texttt{m/44'/60'/0'/0/i}
\end{itemize}

\subsubsection{Seed Phrase (Mnemonic)}

BIP-39 encodes seed as 12-24 words from a 2048-word dictionary.

\begin{examplebox}{Mnemonic Example}
\begin{verbatim}
witch collapse practice feed shame open despair
creek road again ice least
\end{verbatim}

This 12-word phrase represents 128 bits of entropy and can recover the entire wallet.
\end{examplebox}

\begin{securitybox}{Mnemonic Security}
\textbf{Best Practices}:
\begin{itemize}
  \item Never store digitally (no photos, no cloud)
  \item Write on paper, store in secure location
  \item Consider metal backup for fire/water resistance
  \item Never share with anyone
  \item Beware of "recovery" scams asking for your phrase
\end{itemize}
\end{securitybox}

\subsection{Address Generation}

\subsubsection{Bitcoin Address}

\begin{enumerate}
  \item Start with public key (33 or 65 bytes)
  \item Apply SHA-256: $h_1 = \text{SHA256}(pk)$
  \item Apply RIPEMD-160: $h_2 = \text{RIPEMD160}(h_1)$
  \item Add version byte (0x00 for mainnet)
  \item Compute checksum: $c = \text{SHA256}(\text{SHA256}(\text{version} || h_2))[:4]$
  \item Encode in Base58: $\text{Base58}(\text{version} || h_2 || c)$
\end{enumerate}

Result: Address like \texttt{1A1zP1eP5QGefi2DMPTfTL5SLmv7DivfNa}

\subsubsection{Ethereum Address}

\begin{enumerate}
  \item Take public key (64 bytes, uncompressed, without prefix)
  \item Apply Keccak-256 hash
  \item Take last 20 bytes
  \item Add \texttt{0x} prefix
\end{enumerate}

Result: Address like \texttt{0x742d35Cc6634C0532925a3b844Bc9e7595f0bEb}

\subsection{Multi-Signature Wallets}

Require $m$ out of $n$ signatures to authorize transactions.

\textbf{Example}: 2-of-3 multisig
\begin{itemize}
  \item Three key holders (e.g., you, partner, lawyer)
  \item Any two can authorize transaction
  \item Single key compromise doesn't lose funds
\end{itemize}

\textbf{Use Cases}:
\begin{itemize}
  \item Corporate funds (require multiple executives)
  \item Escrow services
  \item Personal backup (key redundancy)
\end{itemize}

\section{Common Security Threats}

\subsection{Private Key Compromise}

\textbf{Attack Vectors}:
\begin{itemize}
  \item Malware/keyloggers
  \item Phishing websites
  \item Clipboard hijacking
  \item Weak random number generators
  \item Physical theft
\end{itemize}

\textbf{Mitigation}:
\begin{itemize}
  \item Use hardware wallets for significant holdings
  \item Verify addresses character-by-character
  \item Use operating system with verified integrity
  \item Never enter private keys on websites
\end{itemize}

\subsection{Transaction Malleability}

Before SegWit, Bitcoin transaction IDs could be modified without invalidating signatures.

\textbf{Problem}: Attackers could change transaction ID, causing wallet confusion.

\textbf{Solution}: Segregated Witness (SegWit) separates signature data, fixing transaction ID calculation.

\subsection{Replay Attacks}

After a blockchain fork, transactions can be valid on both chains.

\textbf{Example}: Bitcoin/Bitcoin Cash split (2017)
\begin{itemize}
  \item Sending BTC on Bitcoin chain also sent BCH on Bitcoin Cash chain
  \item Users lost funds on one chain
\end{itemize}

\textbf{Mitigation}: Replay protection mechanisms, split coins intentionally

\section{Study Questions}

\subsection{Conceptual Questions}

\begin{enumerate}
  \item Explain why reusing the same nonce $k$ in ECDSA is catastrophic.
  \item How does the avalanche effect contribute to blockchain security?
  \item What is the difference between pre-image resistance and collision resistance?
  \item Why can't someone derive a private key from a public key in ECC?
  \item Describe the trade-offs between hot and cold wallet storage.
\end{enumerate}

\subsection{Mathematical Problems}

\begin{enumerate}
  \item If a hash function has 160-bit output, estimate the number of hashes needed to find a collision (use birthday paradox).

  \item Given ECDSA signature $(r, s)$ and two messages $m_1$, $m_2$ signed with the same nonce $k$, derive a formula to recover the private key $sk$.

  \item Calculate the probability that a random private key matches a specific public key for secp256k1 (order $n \approx 2^{256}$).

  \item How many bits of security does SHA-256 provide against collision attacks? Against pre-image attacks?
\end{enumerate}

\subsection{Practical Exercises}

\begin{enumerate}
  \item Implement ECDSA signature generation and verification using a cryptography library.

  \item Create a program to demonstrate the avalanche effect by measuring Hamming distance between hashes of similar inputs.

  \item Write code to generate Bitcoin addresses from public keys (including Base58Check encoding).

  \item Build a simple HD wallet that derives multiple addresses from a single seed.
\end{enumerate}

\section{Further Reading}

\begin{itemize}
  \item NIST FIPS 180-4: Secure Hash Standard (SHA-256 specification)
  \item NIST FIPS 186-4: Digital Signature Standard (ECDSA specification)
  \item Koblitz, N. (1987). \textit{Elliptic curve cryptosystems}. Mathematics of Computation.
  \item BIP-32: Hierarchical Deterministic Wallets
  \item BIP-39: Mnemonic code for generating deterministic keys
  \item Antonopoulos, A. M. (2017). \textit{Mastering Bitcoin}, Chapter 4: Keys, Addresses.
\end{itemize}

\end{document}
